\chapter*{Abstract}
\addcontentsline{toc}{chapter}{Abstract}

This thesis documents an evidence-based engineering adaptation of a multimodal Vision--Language--Action (VLA) driving policy, SimLingo, originally trained from the CARLA autonomous driving simulator to the Quanser QLabs vehicle control platform controlling a QCar~2 vehicle model.

The contribution of this work is a reproducible end-to-end pipeline spanning: (i) expert driving data collection in QLabs, (ii) LoRA-based fine-tuning of an InternVL2-1B backbone using collected expert driving data, and (iii) a real-time inference stack that converts model outputs (waypoint trajectories) into low-level QCar~2 control commands using a simulator-specific PID controller.

The adaptation addresses platform mismatches that materially affect closed-loop behavior, including coordinate frame conventions, control-rate differences, camera resolution preprocessing, and speed estimation consistency with the training data.

Experimental evaluation on a roundabout scenario with five obstacle placement variants shows that the fine-tuned SimLingo model achieves a 60\% overall pass rate (9/15 runs) and a 40\% obstacle avoidance pass rate (4/10 runs), successfully detecting and stopping without contact for obstacles in two of five variants. In two additional variants, the model demonstrates clear obstacle detection and deceleration behavior but makes low-speed bumper contact before fully stopping. Compared against a LiDAR-based adaptive cruise control baseline that achieves 100\% pass rate on the same test scenarios, the vision-only model shows worse but promising performance. Offline validation shows a 25\% reduction in Average Displacement Error (ADE) (from 0.114 to 0.085) over 15~training epochs, converging toward the expert demonstration baseline (ADE${=}0.087$).
